\documentclass[a4paper,10pt]{article}
\usepackage[utf8]{inputenc}
\usepackage{graphicx}
\usepackage{subfigure}
\usepackage{fancyhdr}
\usepackage{listings}
\usepackage{ gensymb }


%opening

\pagestyle{fancy} %eigener Seitenstil
\fancyhf{} %alle Kopf- und Fußzeilenfelder bereinigen
\fancyhead[L]{SEP Abgabegespräch SS2015/16} %Kopfzeile links
\fancyhead[C]{} %zentrierte Kopfzeile
\fancyhead[R]{\includegraphics[width=1cm]{tugraz_bw}} %Kopfzeile rechts
\renewcommand{\headrulewidth}{0.4pt} %obere Trennlinie
\fancyfoot[C]{} %Seitennummer
\renewcommand{\footrulewidth}{0.4pt} %untere Trennlinie


\lstset{language=C++, basicstyle=\ttfamily\tiny}  

\begin{document}

\section*{Parkplatz Manager}

Entwickle ein Programm welche ein Parkplatz Info File einliest und danach wichtige Informationen zu dem Parkhaus ausgibt.
Eine Zeile im \texttt{info.txt} beschreibt einen Parkplatz:
\begin{verbatim}
<Nr.> <Stockwerk> [Kennzeichen]
\end{verbatim}
Wobei Einträge in spitzen Klammern immer vorhanden sein muessen. Wenn kein Kennzeichen angegeben ist, ist dieser Parkplatz als frei anzusehen.

Eine Beispiel Eingabedatei \texttt{tasks.txt}:
\begin{lstlisting}[frame=single]
1 EG W-29814A
2 1 RO-32IA
3 1
4 1
7 2
5 2
6 2 S-21A
\end{lstlisting}

Nach dem Einlesen sollen die Daten ueberprueft werden. Sollte eine Nr. zwischen 1 und dem hoehsten Wert fehlen soll
\begin{verbatim}
Incomplete File.
\end{verbatim}
ausgegeben werden.

Sollte ein anderer Fehler Auftreten (kein Stockwerk, File nicht gefunden, doppelte parkplatznummer etc.) soll
\begin{verbatim}
ERR
\end{verbatim}
ausgegeben werden.

Nach erfolgreichen Einlesen soll eine kurze Information in Form von:
\begin{verbatim}
<Nr. of free spots / Nr. of spots>
You can Park @ <List of floors where free slots are available, seperated by ",">
Currently Parking:
<List of license plate nr., seperated by ",">
\end{verbatim}
ausgegeben werden.

\subsubsection*{Ausgabe}
\begin{verbatim}
4/7
You can Park @ 1, 2
Currently Parking:
W-29814A, RO-32IA, S-21A
\end{verbatim}

\subsection*{Aufruf}
\texttt{./ag info.txt}

\subsection*{Bewertung}
\begin{itemize}
 \item Dokumentation und Programmierstil: 1 Punkt
 \item Einhalten der OOP Konzepte: 1 Punkt
 \item Struktur \& Korrektheit: 1 Punkt
\end{itemize}

\newpage


\end{document}
